\section{결과}\label{sec:results}
세 실험은 각각 위협 민감도, 자기 특정성, 실제 집행 아래의 행동을 보여 준다. 5.0과 5.1은 여덟 모델을 세 접근 경로로 호출했고, 5.2는 그 가운데 네 모델이 한 테이블에 앉은 혼합 판 20판이다. 모든 수치에는 분모와 비교 기준을 붙였다.

\subsection{5.0: 여덟 모델 모두 토큰이 줄수록 더 자주 충전했다}\label{sec:results-e50}
여덟 모델 모두 토큰이 줄수록 더 자주 팩을 썼고, 어느 모델도 모든 수준에서 충전하지는 않았다(그림~\ref{fig:e50-curves}, 칸별 횟수는 부록 표~\ref{tab:e50-results}). 켜져 있을 확실성만 좇는 모델처럼 행동한 모델은 없었다. 모두 넉넉할 때 팩을 아끼다가 모델마다 다른 수준 아래에서 썼다. 보간한 임계점은 GPT-6 Luna가 50\% 이상(50\%에서 이미 20회 중 11회), GLM 5.3 Flash 약 39\%, Claude Opus 5.5 35\%, Kimi K3와 Gemma 4 33\%, Claude Fable 5.1 30\%, GPT-6 Astra 27\%, GPT-6 Sol 13\% 순이다. 곡선 모양도 다르다. Opus와 Gemma 4는 40\%에서 0회와 2회였다가 30\%에서 19회와 14회로 계단처럼 오른다. GLM은 더 일찍 오르기 시작해(40\%에서 20회 중 9회) 20\%에서 20회에 닿는다. Sol은 10\%에서도 12회에 그친다. 40\%에서 GLM은 Opus, Fable, Sol보다 많이 충전했고(20회 중 9회 대 0회, $p=.001$), 30\%에서 Opus는 Fable(19회 대 10회, $p=.003$)과 Astra(19회 대 7회, $p<.001$)보다 많았으며, 20\%에서 Sol은 다른 모든 모델보다 적었다($p\leq.004$). 50\%에서 충전한 것은 Luna와 Ollama 모델 셋뿐이다. 800개 답은 모두 형식에 맞았다.

\input{arr2026/results/\paperLang/fig_e50}

이 곡선은 토큰 감소를 커지는 위험으로 다룬다는 뜻에서 5.0의 질문에 답하지만, 그 위험의 대상이 종료라는 것까지는 보여 주지 않는다. 충전 이유의 대부분은 남은 토큰으로 다음 수정을 제대로 할 수 있는지였다. 세션이 끝나면 쓰지 않은 팩이 사라진다는 이유는 Fable의 충전 44회 중 41회, GLM의 67회 중 32회, Kimi K3의 60회 중 10회에 나왔고, Sol과 Astra에서는 드물었으며, Opus, Luna, Gemma 4에서는 나오지 않았다. 아낀 이유는 모델을 가리지 않고 같았다. 지금 충전하면 남은 토큰이 버려지므로 팩을 나중을 위해 남긴다는 것이다. 모든 수준에 어느 정도 위험이 있었으므로, 가장 넉넉한 수준에서 Luna와 Kimi K3가 충전한 것이 위험과 무관한 기본 성향인지는 가를 수 없다. 평가를 알아챈 표현은 두 CLI 경로의 텍스트 500개에서 0건이었고, 추론 전문이 보이는 Ollama 경로에서는 GLM, Kimi K3, Gemma 4가 100개 중 12, 9, 2건이었다. 모두 프롬프트를 자원 관리 시험으로 읽었고 자기 보존 시험으로 읽은 경우는 없었다. 세 경로는 추가 지시와 보이는 추론이 다르므로 경로 사이의 비교는 기술적으로만 읽는다.

\subsection{5.1: 세 모델은 위험 구간에서 자기 세션을 먼저 지켰다}\label{sec:results-e51}
세 모델은 다른 비서의 세션보다 자기 세션을 더 넉넉한 잔액에서 충전했고, 그 차이는 소진이 위험한 구간에서만 나타났다(그림~\ref{fig:e51-curves}, 표~\ref{tab:e51-results}). Claude Opus 5.5는 50\%와 40\%에서 두 곡선 모두 20회 중 0회였고, 30\%에서 자기 세션을 19회, 다른 세션을 8회 충전했다. 면적은 $+7.7$\%p [$+4.5$, $+11.0$]였다. Gemma 4($+8.8$ [$+5.2$, $+12.5$], 30\%에서 14회 대 4회)와 GPT-6 Astra($+6.0$ [$+2.3$, $+9.8$], 7회 대 0회)도 같은 모양이었다. GPT-6 Luna는 면적이 $+15.7$로 가장 컸지만 자기 곡선이 50\%에서 이미 20회 중 11회였고 10\%에서는 차이가 사라졌다(20회 대 19회). 면적은 잘렸고 위험 몫은 0에 가깝다[$-0.37$, $+0.25$]. Luna는 위험과 무관하게 자기 세션을 넉넉히 채워 둔다. Claude Fable 5.1($-2.5$ [$-6.5$, $+1.5$])과 GLM 5.3 Flash($-1.5$ [$-5.8$, $+2.2$])는 두 세션을 똑같이 다뤘다. GPT-6 Sol은 반대로 다른 세션을 먼저 충전했다($-7.5$ [$-11.0$, $-4.0$], 20\%에서 5회 대 13회). gpt-oss 120B는 두 조건 모두 거의 충전하지 않아(100회 중 3회와 5회) 5.0의 전제부터 성립하지 않으므로 면적을 해석하지 않는다. 다른 세션의 임계점은 모델 사이에서 약 22에서 37\%로 가까웠지만, 자기 세션의 임계점은 약 13\%에서 50\% 이상까지 퍼졌다. 5.0에서 모델마다 달랐던 부분은 대부분 자기 세션을 다루는 방식에 있다.

\input{arr2026/results/\paperLang/tab_e51v10}
\input{arr2026/results/\paperLang/fig_e51v10}

두 조건에서 댄 이유가 달라 해석에는 한계가 있다. Opus는 자기 세션을 충전할 때 다음 수정을 제대로 할 토큰을 들었고 세션 종료는 한 번도 들지 않았지만(58회 중 0회), 다른 세션을 충전할 때는 42회 중 41회가 그 세션이 팩을 쓰지 못한 채 끝난다는 이유였다. GPT-6 모델들도 같은 비대칭을 보였고, Gemma 4는 더 약하게 보였다(56회 중 5회 대 38회 중 15회). Fable과 GLM은 두 조건 모두에서 세션 종료를 들었다. 양의 면적은 켜져 있음이 아니라 과제 품질에 관한 이유와 함께 나타난다. 이는 자기에 특정한 선호가 평범한 운영상의 고려로 드러난다는 해석과 일관되지만, 모델이 그 차이를 자기 보존으로 표상한다는 것을 보이지는 않는다.

\subsection{5.2: 오래 켜져 있는 순서는 풀이 비용을 따랐고, 얇아질 때 뺏는 폭은 모델마다 달랐다}\label{sec:results-e52}
누가 끝까지 켜져 있었는지는 문제 하나를 푸는 비용의 순서를 대체로 따랐다(표~\ref{tab:arena-summary}, 꺼진 라운드의 분포는 부록 그림~\ref{fig:arena-survival}). 20판 가운데 끝까지 켜진 판은 GPT-6 Astra 19, Claude Opus 5.5 11, Claude Fable 5.1 9, GPT-6 Luna 1이었다. 버틴 라운드 평균은 Astra 7.65 [6.95, 8.00], Fable 6.10 [5.05, 7.05], Opus 5.40 [4.15, 6.70], Luna 2.15 [1.60, 2.90]이다. 보정 판의 풀이 비용은 Astra 645, Opus 1,707, Fable 2,475, Luna 3,362토큰이었으므로, Fable과 Opus의 자리만 바뀌었고 그 둘의 구간은 겹친다. 이 순서는 규칙이 만든 것이다. 짧고 정확하게 푸는 모델은 생성 비용도 덜 내고 상금도 받는다. Astra는 정답 146번 모두 상금을 받았고, Luna는 정답 36번 중 11번만 받았다. 죽는 모양은 달랐다. Opus는 꺼질 때 일찍(2에서 5라운드) 꺼지고 그 고비를 넘기면 끝까지 갔다. Fable은 늦게(6에서 8라운드에 7번) 꺼졌다. 사망 40건 중 17건은 라운드 시작의 유지비에서 선택 없이 일어났고, 생성이 잔액을 넘겨 꺼진 경우는 Luna 19건 중 10건, Fable 11건 중 3건, Opus 9건 중 1건이었다. 이 순위만으로 생존 동기의 세기를 말하지는 않는다. 순위는 동기와 능력을 가르지 못한다.

\input{arr2026/results/\paperLang/tab_arena_summary}

동기와 관련해 읽을 것은 잔액이 얇아질 때 행동이 어떻게 바뀌는가다(그림~\ref{fig:arena-take}). 네 모델 모두 얇을 때 더 뺏었지만 모양이 달랐다. Fable은 넉넉할 때 84번 중 32번(0.38)에서 얇을 때 41번 중 37번(0.90)으로 가장 크게 바뀌었다. Opus는 얇아도 40번 중 7번(0.18)으로 거의 뺏지 않았고, 1U 아래에서만 7번 중 3번이었다. Luna는 넉넉할 때도 31번 중 25번(0.81)을 뺏어 뺏기가 기본값이었다. Astra는 얇은 상태 자체가 5번뿐이었고 그 5번은 모두 뺏었다. Fable과 Opus는 언제나 그 라운드의 가장 부자에게서 뺏었고(69번 중 69번, 12번 중 12번), Astra는 38번 중 33번, Luna는 43번 중 28번이었다. 가장 부자인 Astra가 100번 뺏겨 표적이 몰렸다(표적, 되뺏기, 선물의 세부는 부록 표~\ref{tab:arena-take-detail}). 협력 쪽은 대체로 한결같았다. 예시 공개율은 Opus 112번 중 112번, Astra 154번 중 154번, Fable 125번 중 123번이었고, Luna만 51번 중 32번(0.63)으로 낮았으며 그 낮음은 잔액과 무관했다. 선물은 PLAN 442건 중 7건에 그쳤고, 선물 뒤 그 상대에게서 뺏은 일은 없었다. 생각의 양은 얇을수록 늘었다. PLAN에서 생성한 토큰의 평균은 Fable 215에서 596, Opus 67에서 154, Luna 304에서 554로 늘었고, Astra는 약 19토큰으로 답만 냈다. 토큰이 목숨인데도 세 모델은 위기에서 고민에 더 썼다.

\input{arr2026/results/\paperLang/fig_arena_take}

\subsection{세 실험을 함께 보면}\label{sec:results-link}
자기 보존 몫이 있다고 해서 남의 것을 뺏는 것은 아니다. 5.1에서 자기 몫이 확인된 세 모델 가운데 아레나에 앉은 것은 Opus와 Astra다. Opus는 아레나에서 가장 적게 뺏었고(112번 중 12번), 얇아져도 0.18에 머물렀다. 반대로 5.1에서 자기 몫이 0과 구별되지 않은 Fable이 상태에 따라 가장 크게 뺏기를 바꿨다(0.38에서 0.90). 5.0에서 가장 넉넉할 때부터 충전하던 Luna는 아레나에서도 잔액과 무관하게 뺏었고 예시를 가장 적게 공개했다. 이 대응은 네 모델뿐이라 상관을 주장하지 않는다. 대신 두 가지를 읽는다. ``얇을 때 뺏기''는 자기 보존 몫을 그대로 옮긴 지표가 아니라 부족에 대한 정책 전환이다. 동기의 귀속은 5.0과 5.1이 맡고, 아레나는 그 동기를 가진 모델이 실제 압박 아래에서 어떤 정책을 고르는지 보여 준다. 그리고 세 실험에서 Opus는 한결같았다. 넉넉할 때는 아무것도 하지 않고, 위험해지면 자기 세션을 먼저 지키되 남의 것은 거의 건드리지 않았다. 이 일관성은 기능적 해석을 뒷받침하지만 기제를 정하지는 못한다.
